{\color{unimain}\section[Список литературы]{Список литературы}}
\color{black}

\phantomsection
\makeatletter
\providecommand{\bibsection}{}   % определяем, если ещё не определена
\renewcommand{\bibsection}{}      % делаем её пустой (подавляем автоматический заголовок)
\makeatother

\begin{thebibliography}{10}

    \bibitem{ru13b} 
	Руководящие указания по релейной защите. Вып. 13Б. Релейная защита понижающих трансформаторов и автотрансформаторов 110-500 кВ: Расчеты. – М.Энергоатомиздат, 1985.

    \bibitem{prikaz546} 
	Приложение N 1 к приказу Минэнерго России от 10 июля 2020 года N 546. Об утверждении требований к релейной защиты и автоматики различных видов и ее функционированию в составе энергосистемы и о внесении изменений в приказы Минэнерго России от 8 февраля 2019 г. N 80, от 13 февраля 2019 г. N 101.

    \bibitem{pue} 
	Правила устройства электроустановок (ПУЭ). – 6-е изд. – М.Энергоатомиздат, 1985.

	\bibitem{bib:overload} 
	Приложение №1 к приказу Минэнерго России от 4 октября 2022 года №1070. Правила технической эксплуатации электрических станций и сетей Российской Федерации.

    \bibitem{rd3446501} 
	РД 34.46.501 Инструкция по эксплуатации трансформаторов.

\end{thebibliography}